\section{Problem formulation}
\label{sec:problem}

\subsection{Production and logistics}
\label{sec:problem:model}

We study the production--logistics collaborative scheduling problem (PLCSP).
The problem couples two decision layers that are usually solved apart:
a flexible job shop, where each operation must be assigned to one eligible
machine, and a fleet of automated guided vehicles (AGVs),
which must move every job between machines and to and from a load--unload
station. Assigning an operation to a machine is a decision; the order in which
a machine works through the jobs waiting at it is not one, and follows the
order in which the vehicles deliver them. The term follows the recent
literature on production--logistics
collaborative scheduling; the same problem class appears in the older
literature as integrated scheduling of machines and AGVs.

Let $\mathcal{J}$ be the set of $n$ jobs, $\mathcal{M}$ the set of $m$
machines, and $\mathcal{V}$ the set of $V$ vehicles. Job $j$ consists of
$n_j$ operations that must be processed in the fixed order
$O_{j,1} \to O_{j,2} \to \dots \to O_{j,n_j}$. Operation $O_{j,k}$ may be
processed on any machine in the eligible set $\mathcal{M}_{j,k} \subseteq
\mathcal{M}$; its processing time on machine $i$ is $p_{jki}$. If job
$j'$ is scheduled on a machine immediately after job $j$, a sequence-dependent
setup time $s_{jj'}$ elapses before $j'$ can start.

\paragraph{Shop layout and transport}
Machines occupy the cells of a two-dimensional grid with $n_{\mathrm{col}} =
\lceil \sqrt{m} \rceil$ columns and $n_{\mathrm{row}} = \lceil m /
n_{\mathrm{col}} \rceil$ rows. Empty cells are permitted. The aisles are the
grid lines between cells, so a passage never crosses a machine. The aisle
network is a graph whose nodes are the $(n_{\mathrm{row}}+1)(n_{\mathrm{col}}+1)$
intersections and whose edges are the horizontal and vertical segments between
adjacent intersections; the weight of an edge is its physical length. Each
machine has a dock, connected to the nearest corner intersection of its cell.
The load--unload (LU) station is a node outside the grid, joined to the left
boundary by one connecting segment. Jobs enter the shop at the LU station and
return there when their last operation is complete. Completion times are
therefore measured at the LU station, so the makespan of a job includes both
its entry travel and its return travel. Figure~\ref{fig:layout} shows the
layout.

\begin{figure}[htbp]
\centering
\includegraphics[width=\columnwidth]{fig_layout.pdf}
\caption{Shop layout. Machines occupy grid cells; the aisles are the grid lines
between cells, so a passage never crosses a machine. Each machine has a dock
linking it to the nearest intersection. The load--unload station lies outside
the grid and is joined to the left boundary by one connecting segment. Charging
stations occupy aisle intersections. The solid line is one AGV travel.}
\label{fig:layout}
\end{figure}

Two transport-time calibers coexist in this study, and the caliber is a
property of the instance rather than of the run configuration. Under the
geometric caliber, the travel time between two nodes is the shortest-path
length in the aisle graph divided by the effective vehicle speed. Under the
matrix caliber, the travel time is read from a machine-to-machine travel-time
matrix supplied with the instance; entries are already in minutes, and the
matrix includes the LU station as its zeroth row and column. The matrix is
asymmetric, so a lookup preserves direction. A fleet-wide speed multiplier
still applies under the matrix caliber; the aisle width and the nominal
vehicle speed do not.

\paragraph{Constraints}
Ten constraints are modelled. Each is an independent switch, so a run can
enable any subset.

\begin{description}\itemsep2pt
\item \textbf{C1 --- AGV congestion, zone conflict, deadlock.}
      An aisle segment carries at most one vehicle. A vehicle requests the next
      segment before releasing the previous one, and wait cycles are detected
      and broken.
\item \textbf{C2 --- Finite input and output buffers.}
      A machine blocks when its output buffer is full, and an operation waits
      when the input buffer of its machine is full.
\item \textbf{C3 --- Machine failure.}
      Failures occur at a rate that rises with machine age, and repair occupies
      the machine for a random interval.
\item \textbf{C4 --- Rework.}
      A completed operation is re-queued with probability $p_{\mathrm{rw}}$.
\item \textbf{C5 --- Sequence-dependent setup.}
      A setup time $s_{jj'}$ precedes each job change on a machine.
\item \textbf{C6 --- Due dates and tardiness.}
      Per-job due dates enter Equation~\eqref{eq:twt}. This is the one
      constraint that changes only what is measured: the simulation trajectory
      is unaffected.
\item \textbf{C7 --- AGV failure.}
      A vehicle stops for a repair interval, and tasks already assigned to it
      are returned to the dispatcher and re-offered to the remaining vehicles.
\item \textbf{C8 --- Heterogeneous fleet.}
      Vehicles differ in speed multiplier, load capacity, and battery size.
\item \textbf{C9 --- Charging and battery.}
      Vehicle energy is tracked. A depleted vehicle becomes unavailable, and
      charging is a scheduled decision.
\item \textbf{C10 --- Preventive maintenance.}
      Each machine carries a maintenance clock, and maintenance is a scheduled
      decision that resets it.
\end{description}

The distinction between a constraint that changes the dynamics and one that
changes only the objectives matters for the ablations in
Section~\ref{sec:results}. Switching off a dynamics constraint changes the
trajectory; switching off C6 leaves every trajectory bit-identical and only
removes a term from the objective.

\subsection{Objectives}
\label{sec:problem:objectives}

Three objectives are reported. Let $C_j$ denote the completion time of job
$j$, measured when the job reaches the LU station, and let $d_j$ denote its
due date.

\paragraph{Makespan}
\begin{equation}
\Cmax = \max_{j \in \mathcal{J}} C_j .
\label{eq:makespan}
\end{equation}

\paragraph{Total energy consumption}
Energy is accounted by a three-state machine model and a three-state vehicle
model, plus a fixed shop load. Let $\pi^{\mathrm{p}}_i$, $\pi^{\mathrm{s}}_i$,
and $\pi^{\mathrm{u}}_i$ be the processing, setup, and standby power of machine
$i$, in kilowatts, and let $T^{\mathrm{p}}_i$, $T^{\mathrm{s}}_i$, and
$T^{\mathrm{u}}_i$ be the total minutes machine $i$ spends in each state. Let
$\pi^{\mathrm{l}}$, $\pi^{\mathrm{e}}$, and $\pi^{\mathrm{i}}$ be the loaded,
empty-travel, and standby power of a vehicle, and let $T^{\mathrm{l}}$,
$T^{\mathrm{e}}$, and $T^{\mathrm{i}}$ be the corresponding fleet totals. Let
$P_0$ be the fixed shop power. Then
\begin{equation}
E_{\mathrm{mach}} = \frac{1}{60} \sum_{i \in \mathcal{M}}
   \left( \pi^{\mathrm{p}}_i T^{\mathrm{p}}_i
        + \pi^{\mathrm{s}}_i T^{\mathrm{s}}_i
        + \pi^{\mathrm{u}}_i T^{\mathrm{u}}_i \right),
\label{eq:energy:mach}
\end{equation}
\begin{equation}
E_{\mathrm{agv}} = \frac{1}{60} \left( \pi^{\mathrm{l}} T^{\mathrm{l}}
   + \pi^{\mathrm{e}} T^{\mathrm{e}} + \pi^{\mathrm{i}} T^{\mathrm{i}} \right),
\label{eq:energy:agv}
\end{equation}
\begin{equation}
E_{\mathrm{tot}} = E_{\mathrm{mach}} + E_{\mathrm{agv}}
   + \frac{P_0 \, \Cmax}{60} .
\label{eq:energy}
\end{equation}
The fixed shop load is charged over the makespan, so the third term is linear
in $\Cmax$ by construction.

Section~\ref{sec:results} also reports a second energy account,
$E_{\mathrm{net}}$, which keeps only the terms attributable to the work
actually performed:
\begin{equation}
E_{\mathrm{net}} = \frac{1}{60} \sum_{i \in \mathcal{M}}
   \left( \pi^{\mathrm{p}}_i T^{\mathrm{p}}_i
        + \pi^{\mathrm{s}}_i T^{\mathrm{s}}_i \right)
   + \frac{1}{60} \left( \pi^{\mathrm{l}} T^{\mathrm{l}}
        + \pi^{\mathrm{e}} T^{\mathrm{e}} \right) .
\label{eq:energy:net}
\end{equation}
$E_{\mathrm{net}}$ excludes machine standby, vehicle standby, and the fixed
shop load. The two accounts behave differently as objectives, and
Section~\ref{sec:results} reports that difference as a result rather than
treating the choice as a formality.

\paragraph{Total weighted tardiness}
\begin{equation}
\TWT = \sum_{j \in \mathcal{J}} \wj \max(0, \, C_j - d_j).
\label{eq:twt}
\end{equation}
All job weights are set to one in this study, so \TWT\ is the unweighted total
tardiness. The weights are retained in the formulation because the objective
is defined with them and because the instance data carry a weight field.

\subsection{Due dates}
\label{sec:problem:due}

Due dates must be exogenous, per job, and spread. A common due date makes the
tardiness objective degenerate: if all $d_j$ are equal, then $\TWT = 0$
whenever $\Cmax \le d$, so the objective carries no signal beyond the
makespan. Two formulations that are natural at first sight fail on this point,
and the failure is structural rather than numerical. Setting $d_j$ to a fixed
multiple of the job's own processing time ignores queueing and transport, so
every job is tardy. Setting $d_j = \tau M_{\mathrm{ref}}$, where
$M_{\mathrm{ref}}$ is the makespan of a reference schedule on the same
instance, is endogenous: changing the fleet size changes
$M_{\mathrm{ref}}$ and therefore changes the due dates.

We therefore generate due dates from the job work content. Let
$W_j = \sum_{k} \min_{i \in \mathcal{M}_{j,k}} p_{jki}$ be the shortest
attainable processing time of job $j$, and let
\begin{equation}
\mathrm{LB} = \max \left( \frac{1}{m} \sum_{j \in \mathcal{J}} W_j, \;
   \max_{j \in \mathcal{J}} W_j \right)
\label{eq:lb}
\end{equation}
be a load lower bound on the makespan of any feasible schedule. Let
$\rho_j \in [0,1]$ be the rank of $W_j$ among all work contents, normalized
to the unit interval, with $\rho_j = 0.5$ when $n = 1$. The due date of job
$j$ is
\begin{equation}
d_j = \mathrm{LB} \cdot \tau \cdot \left( 1 + R (2\rho_j - 1) \right).
\label{eq:duedate}
\end{equation}
Equation~\eqref{eq:duedate} reads only instance data: the reference schedule
and the fleet configuration do not appear, so changing the fleet size does not
move any due date. The pair $(\tau, R)$ is calibrated per instance. The
tardiness rate of the reference schedule and the tardiness rate after a 12\%
improvement are both required to fall in $[0.20, 0.75]$; among the pairs that
satisfy this, the calibration selects the one minimizing
$|\eta_{\mathrm{ref}} - 0.45| + |\eta_{\mathrm{imp}} - 0.40| + 0.02R$, where
$\eta$ denotes a tardiness rate. Tightening $\tau$ only increases our own
tardiness, so the rule carries no incentive to tune the target in our favour.
$R$ is bounded below by $0.2$ because $R = 0$ collapses
Equation~\eqref{eq:duedate} to a common due date, which is the degeneracy this
section exists to avoid. It is bounded above by $0.8$ because the factor
$1 + R(2\rho_j - 1)$ crosses zero once $R > 1$, which would produce negative
due dates.

The calibrated $\tau$ differs across instances by up to a factor of 3.4 under
the geometric caliber and 12.0 under the matrix caliber. That spread is
inherent to work-content-based due dates and is reported rather than hidden. Table~\ref{tab:duedates} gives the frozen pairs for both
transport calibers. They are part of the instance data in the sense that
matters here --- the reward of Equation~\eqref{eq:reward} is not reproducible
without them --- so they are reported in full even though the policy results in
this paper come from a single instance.

Both columns were produced by a search over a fixed grid, with a step of 0.05
in $\tau$ and 0.10 in $R$, rather than chosen by hand. Halving the grid step
leaves the induced tardiness rates unchanged on nine of the ten instances under
the geometric caliber; on the tenth it moves to a neighbouring pair with a
slightly worse score. The table reports the frozen values.

\begin{table}[htbp]
\centering
\caption{Frozen due-date parameters, one pair per caliber. The two calibers
measure travel time on different scales, so they need separate pairs.}
\label{tab:duedates}
\setlength{\tabcolsep}{4pt}
\begin{threeparttable}
\begin{tabular}{@{}lrrrr@{}}
\toprule
 & \multicolumn{2}{c}{Geometric} & \multicolumn{2}{c}{Matrix} \\
\cmidrule(r){2-3}\cmidrule(l){4-5}
Instance & $\tau$ & $R$ & $\tau$ & $R$ \\
\midrule
mk01 & 2.75 & 0.3 & 19.80 & 0.4 \\
mk02 & 2.45 & 0.4 & 22.15 & 0.2 \\
mk03 & 3.65 & 0.4 & 16.10 & 0.3 \\
mk04 & 5.00 & 0.7 & 20.45 & 0.5 \\
mk05 & 1.70 & 0.5 &  4.40 & 0.4 \\
mk06 & 5.85 & 0.6 & 45.50 & 0.5 \\
mk07 & 1.90 & 0.5 &  8.60 & 0.3 \\
mk08 & 2.70 & 0.8 &  8.80 & 0.6 \\
mk09 & 2.70 & 0.8 & 13.00 & 0.4 \\
mk10 & 3.10 & 0.8 & 52.70 & 0.6 \\
\bottomrule
\end{tabular}
\begin{tablenotes}\footnotesize
\item Frozen values; the grid search that produced them is described in the
      text.
\end{tablenotes}
\end{threeparttable}
\end{table}

The tardiness objective also saturates. Because due dates are absolute
thresholds, a schedule that finishes every job before the earliest $d_j$
attains $\TWT = 0$ regardless of how much better it is. The calibration above
pins the reference schedule near the middle of the band rather than at either
edge, but the saturation remains, and it is reported wherever it occurs rather
than being treated as a measurement of zero tardiness.
